% ===========================================================================
% Chapitre 3 — Architecture et conception.
% Toute la matière de ce chapitre est relevée dans le dépôt au 2026-09-22.
% ===========================================================================
\chapter{Architecture et conception}
\label{chap:architecture}

\section{Introduction}
\label{sec:arch_intro}

Ce chapitre traduit les exigences du chapitre précédent en une architecture.
Il énonce d'abord les objectifs architecturaux et les réponses apportées,
délimite le système et ses dépendances, décrit le découpage en composants et
le modèle de données, expose les séquences des deux parcours les plus
structurants, analyse les frontières de confiance, décrit le déploiement, et
consigne enfin les décisions d'architecture avec leurs alternatives écartées.

\section{Objectifs architecturaux}
\label{sec:arch_objectives}

Le tableau~\ref{tab:arch_objectives} associe à chaque objectif architectural
la réponse retenue.

\begin{table}[htbp]
  \centering
  \caption{Objectifs architecturaux et réponses apportées}
  \label{tab:arch_objectives}
  \begin{tabularx}{\textwidth}{L{3.6cm} X}
    \toprule
    \textbf{Objectif} & \textbf{Réponse architecturale} \\
    \midrule
    Explicabilité      & Le score est stocké décomposé en sous-scores persistés, non recalculé à l'affichage : ce qui a été montré au recruteur reste consultable ensuite. \\
    Réactivité perçue  & Les traitements longs — analyse de CV, indexation — s'exécutent hors du cycle de la requête, dans une session de base dédiée. \\
    Résilience         & Toute dépendance externe est facultative et dotée d'un repli implémenté ; aucune n'est sur le chemin critique. \\
    Sécurité           & Politique d'autorisation externalisée, refus par défaut, et traitement des liens publics à jeton comme une frontière de confiance à part entière. \\
    Intégrité du schéma& Le schéma est piloté exclusivement par migrations ; l'application vérifie la révision au démarrage et refuse de servir un schéma désynchronisé. \\
    Confidentialité    & Aucune donnée personnelle de CV n'est écrite dans les journaux, à aucun niveau de verbosité. \\
    \bottomrule
  \end{tabularx}
\end{table}

\section{Contexte système}
\label{sec:system_context}

La figure~\ref{fig:system_context} délimite le système, ses acteurs et ses
dépendances externes. Le tableau~\ref{tab:external_deps} détaille chaque
dépendance et son repli.

\pqfigure{system_context}
  {Frontière du système PIQBIT, acteurs et dépendances externes. Le trait
   discontinu marque une dépendance \emph{facultative} dotée d'un repli
   implémenté.}
  {fig:system_context}

\begin{table}[htbp]
  \centering
  \caption{Dépendances externes et comportement de repli}
  \label{tab:external_deps}
  \begin{tabularx}{\textwidth}{L{3.2cm} L{4.2cm} X}
    \toprule
    \textbf{Dépendance} & \textbf{Usage} & \textbf{Comportement si indisponible} \\
    \midrule
    Serveur SMTP     & Invitations, confirmations, rappels & Le message est journalisé au lieu d'être émis ; le parcours se poursuit \\
    Google Calendar  & Synchronisation des entretiens confirmés & Un fichier \texttt{.ics} téléchargeable remplace la synchronisation \\
    Google Meet      & Lien de visioconférence & \textbf{Aucun repli} : sans compte Google rattaché au recruteur désigné, aucun lien n'est créé et aucun rappel de connexion n'est envoyé \\
    Google OAuth     & Connexion par identité externe & La connexion par mot de passe reste disponible \\
    Modèles de langage & Enrichissement de l'assistant & L'assistant reste extractif : il cite la base, sans reformulation \\
    Dépôt de modèles & Téléchargement du modèle de plongement & Le cache local est employé ; à défaut, la recherche sémantique est désactivée sans bloquer la notation \\
    \bottomrule
  \end{tabularx}
\end{table}

Cette propriété a un coût qu'il faut énoncer : des chemins de repli moins
éprouvés que le chemin nominal. La campagne de validation n'a ainsi exercé que
le repli de la recherche sémantique, et non son chemin vectoriel natif
(section~\ref{sec:validation_results}).

\section{Architecture en composants}
\label{sec:component_view}

La figure~\ref{fig:component_view} présente le découpage en composants, et le
tableau~\ref{tab:components} les responsabilités de chaque couche.

\pqfigure{component_view}
  {Vue logique en composants : client Angular, frontière HTTP, couche API,
   services applicatifs et persistance. Relevé du 22 septembre 2026.}
  {fig:component_view}

\begin{table}[htbp]
  \centering
  \caption{Composants de PIQBIT et responsabilités}
  \label{tab:components}
  \begin{tabularx}{\textwidth}{L{3.4cm} X}
    \toprule
    \textbf{Composant} & \textbf{Responsabilité} \\
    \midrule
    Client Angular       & Interfaces front-office, back-office et pages publiques ; chargement différé des modules lourds \\
    Serveur statique     & Service du bundle, relais des appels d'interface et des liaisons temps réel, en-têtes de sécurité \\
    Couche API           & 18 modules d'endpoints exposant 173 routes et deux points de liaison temps réel \\
    Authentification     & Vérification du jeton, second facteur, identité externe \\
    Autorisation         & Décision d'accès par rôle, refus par défaut \\
    Validation           & Contrôle des entrées aux frontières par schémas typés \\
    Services applicatifs & 33 modules portant les règles métier de chaque domaine \\
    Intégrations         & Trois clients vers les services externes, tous facultatifs \\
    Accès aux données    & Correspondance objet-relationnel asynchrone sur 25 tables \\
    \bottomrule
  \end{tabularx}
\end{table}

La règle de dépendance est unidirectionnelle : l'interface appelle l'API, qui
appelle les services, qui appellent la persistance. Un service n'appelle jamais
un endpoint, et la persistance ne connaît aucune couche supérieure. Cette
discipline rend chaque couche testable indépendamment.

\section{Couche applicative et interfaces}
\label{sec:interfaces}

Le tableau~\ref{tab:backend_modules} récapitule les modules d'endpoints et
leur rôle architectural.

\begin{table}[htbp]
  \centering
  \caption{Modules d'endpoints et rôles architecturaux}
  \label{tab:backend_modules}
  \begin{tabularx}{\textwidth}{L{3.4cm} X}
    \toprule
    \textbf{Module} & \textbf{Rôle} \\
    \midrule
    Authentification  & Inscription, connexion, second facteur, réinitialisation, identité externe \\
    Comptes           & Profils, rôles, photo, journal d'audit, suppression de compte \\
    Recrutement       & Offres, candidatures, dépôt et gestion des CV \\
    Notation          & Analyse de CV, scores, classement \\
    Sémantique        & Indexation vectorielle et recherche par proximité de sens \\
    Assistant         & Interrogation de la base de connaissances, réponses sourcées \\
    Entretiens        & Invitations, créneaux, confirmation publique, agenda, visioconférence \\
    Évaluation        & Questions, sessions de test, correction, résultats \\
    Négociation       & Tours d'échange, assistance à la contre-proposition \\
    Contrats          & Édition, envoi, signature publique, certificat \\
    Employés          & Répertoire et fiches employé \\
    Budget            & Enveloppes par département et imputation des dépenses \\
    Carrière          & Objectifs et jalons de progression \\
    Rapports          & Synthèses et instantanés archivés \\
    Décision          & Comparaison de candidats et aide au choix \\
    Notifications     & Diffusion en temps réel et historique \\
    Tableau de bord   & Indicateurs agrégés du back-office \\
    Flux de travail   & Transitions de statut des candidatures \\
    \bottomrule
  \end{tabularx}
\end{table}

Les conventions d'interface sont uniformes : validation des entrées par schémas
typés~\cite{pydantic}, codes de statut normalisés, et identifiants universels comme clés
externes. Deux points de liaison temps réel complètent les routes : l'un pour
les notifications, l'autre pour la négociation, tous deux authentifiés par
jeton — une liaison persistante ne portant pas les en-têtes d'une requête
ordinaire.

\section{Modèle de données et persistance}
\label{sec:data_model}

Le modèle comporte 25 tables, groupées par domaine
(figure~\ref{fig:domain_model}). Le tableau~\ref{tab:domain_groups} en résume
les responsabilités.

\pqfigurewide{domain_model}
  {Modèle de données : 25 tables groupées par domaine. Les cardinalités
   proviennent des contraintes réelles ; deux liens non contraints par le
   moteur sont signalés.}
  {fig:domain_model}

\begin{table}[htbp]
  \centering
  \caption{Groupes du domaine et responsabilités de persistance}
  \label{tab:domain_groups}
  \begin{tabularx}{\textwidth}{L{3.4cm} L{2.0cm} X}
    \toprule
    \textbf{Groupe} & \textbf{Tables} & \textbf{Responsabilité} \\
    \midrule
    Identité et accès    & 2 & Comptes, rôles, second facteur, journal d'audit \\
    Recrutement          & 4 & Candidats, offres, candidatures, favoris \\
    Analyse et notation  & 4 & Analyses de CV, scores, vecteurs et segments indexés \\
    Évaluation           & 5 & Invitations, créneaux, questions, sessions et réponses de test \\
    Embauche             & 4 & Négociations, tours, contrats, employés \\
    RH interne et diffusion & 4 & Budgets, dépenses, plans de carrière, notifications \\
    Instantanés          & 2 & Rapports archivés et bases de connaissances \\
    \bottomrule
  \end{tabularx}
\end{table}

\paragraph{Décision D-01 — le schéma est piloté par migrations.}
Les révisions sont gérées par Alembic~\cite{alembic}, l'outil de migration
associé à la couche de persistance~\cite{sqlalchemy}.
Le schéma n'est jamais créé ni modifié par l'application. Au démarrage,
celle-ci \emph{vérifie} que la base porte la dernière révision de migration et
journalise une erreur explicite en cas d'écart. L'alternative écartée — la
création automatique des tables au démarrage — produit une dérive silencieuse
entre environnements : une colonne ajoutée en développement apparaît sans
migration, et son absence en production ne se découvre qu'au premier appel qui
la touche. Le critère retenu est l'absence de dérive, garantie par un contrôle
en intégration continue qui compare les migrations aux modèles.

\section{Séquences principales}
\label{sec:dynamic_views}

Deux parcours portent l'essentiel de la dynamique du système.

\subsection{Candidature, analyse et notation}

La figure~\ref{fig:sequence_scoring} détaille le parcours du dépôt d'une
candidature jusqu'au classement du candidat.

\pqfigure{sequence_scoring}
  {Candidature, analyse de CV, notation et classement. La réponse
   \texttt{201} part avant le traitement : l'analyse s'exécute en tâche de
   fond, dans une session de base dédiée.}
  {fig:sequence_scoring}

Le point de conception à retenir est la frontière asynchrone. La réponse est
émise dès la candidature enregistrée, avant l'analyse — qui prend plusieurs
secondes. Cela impose deux contraintes : la tâche de fond doit ouvrir sa propre
session de base, celle de la requête étant déjà fermée ; et l'interface doit
représenter un état transitoire, l'analyse n'étant pas immédiatement
disponible. Le bénéfice est une latence perçue constante, indépendante de la
taille du CV.

\subsection{Contrat, signature et activation de l'employé}

La figure~\ref{fig:sequence_contract} détaille la signature d'un contrat par
lien public.

\pqfigure{sequence_contract_signature}
  {Contrat, signature par lien à jeton et activation de l'employé. Le jeton
   tient lieu d'authentification : les gardes d'état précèdent toute écriture.}
  {fig:sequence_contract}

Deux propriétés sont conçues explicitement. D'une part les \emph{gardes
d'état} : un contrat déjà signé, refusé, ou dans un état autre qu'envoyé, est
rejeté avec un code distinct selon le cas — un lien rejoué ne signe jamais deux
fois. D'autre part l'\emph{idempotence} de la création d'employé : la fiche est
recherchée par adresse avant d'être créée, ce qui évite les doublons si
l'opération est rejouée.

\section{Sécurité et frontières de confiance}
\label{sec:security_architecture}

La figure~\ref{fig:security_view} distingue les deux chemins d'entrée du
système, et le tableau~\ref{tab:security_controls} récapitule les contrôles.
Les frontières retenues suivent les catégories de risque du référentiel
OWASP~\cite{owasptop10}, en particulier le contrôle d'accès défaillant et
l'exposition de données sensibles.

\pqfigure{security_view}
  {Frontières de confiance. Deux chemins d'entrée coexistent : authentifié
   (jeton signé, second facteur, autorisation) et public à jeton — neuf routes
   où le jeton tient lieu d'authentification.}
  {fig:security_view}

\begin{table}[htbp]
  \centering
  \caption{Frontières de confiance et contrôles associés}
  \label{tab:security_controls}
  \begin{tabularx}{\textwidth}{L{3.4cm} X}
    \toprule
    \textbf{Frontière} & \textbf{Contrôle} \\
    \midrule
    Entrée réseau        & En-têtes de sécurité, origine unique. \textbf{Aucune terminaison TLS} dans la composition fournie \\
    Authentification     & Jeton signé à durée limitée ; second facteur par mot de passe à usage unique \\
    Autorisation         & Politique déclarative externalisée, refus par défaut, héritage de rôles \\
    Liens publics à jeton& Jeton de 32 octets imprévisible, gardes d'état avant écriture, durée de validité bornée \\
    Dépôt de fichiers    & Liste blanche de types et plafond de taille appliqués avant tout traitement \\
    Journalisation       & Aucune donnée personnelle de CV, à aucun niveau de verbosité \\
    Piste d'audit        & Acteur, action et cible consignés pour les opérations sensibles \\
    \bottomrule
  \end{tabularx}
\end{table}

\paragraph{Décision D-02 — la politique d'autorisation est externalisée.}
Les droits sont décrits dans un fichier déclaratif, et non dispersés en
conditions dans le code. La politique parle du métier — « le recruteur peut
lire les candidatures » — plutôt que des routes ; chaque endpoint déclare la
ressource et l'action qu'il engage. L'alternative écartée — un contrôle écrit
sur les chemins et les verbes — se désynchronise au premier renommage de route
et se relit mal. Le critère retenu est l'auditabilité : la politique complète
se lit dans un seul fichier, ce que la section~\ref{sec:failure_analysis}
montre justement utile lorsqu'un écart survient.

\paragraph{Les liens publics à jeton sont une frontière à part entière.}
Neuf routes acceptent un jeton en guise d'authentification : confirmation
d'entretien, passage de test, signature et refus de contrat. C'est un choix
fonctionnel — exiger un compte d'un candidat pour confirmer un créneau
ajouterait une friction disproportionnée — mais il déplace la sécurité du
porteur vers le jeton. D'où trois exigences de conception : entropie suffisante
(32 octets), validité bornée dans le temps, et gardes d'état systématiques
avant toute écriture.

\section{Vue de déploiement}
\label{sec:deployment_view}

La figure~\ref{fig:deployment_view} restitue la composition de
conteneurs~\cite{docker} au 22 septembre 2026.

\pqfigure{deployment_view}
  {Vue de déploiement. Le conteneur applicatif ne publie aucun port : il n'est
   atteignable qu'à travers le serveur statique. Aucune terminaison TLS, aucun
   environnement cible.}
  {fig:deployment_view}

Trois constats doivent être énoncés sans détour. Le conteneur applicatif
\emph{ne publie aucun port} vers l'hôte : il n'est joignable qu'à travers le
serveur statique, et c'est une propriété voulue. En revanche, la composition
\emph{ne comporte aucune terminaison TLS} : le point d'entrée est en clair. Et
\emph{aucun environnement d'hébergement cible} n'existe à ce jour : la
composition décrite est celle du poste de développement, sans procédure de
retour arrière définie.

\section{Décisions d'architecture}
\label{sec:architecture_decisions}

Le tableau~\ref{tab:arch_decisions} récapitule les décisions structurantes.
Chacune énonce l'alternative écartée et le critère d'arbitrage, faute de quoi
une décision n'est qu'une préférence.

\begin{table}[htbp]
  \centering
  \caption{Décisions d'architecture principales}
  \label{tab:arch_decisions}
  \begin{tabularx}{\textwidth}{L{1.3cm} L{3.4cm} X}
    \toprule
    \textbf{Réf.} & \textbf{Décision} & \textbf{Alternative écartée et critère} \\
    \midrule
    D-01 & Schéma piloté par migrations, vérifié au démarrage
         & \emph{Écarté :} création automatique des tables. \emph{Critère :} absence de dérive entre environnements \\
    D-02 & Politique d'autorisation externalisée, refus par défaut
         & \emph{Écarté :} contrôles dispersés par route. \emph{Critère :} auditabilité en un seul fichier \\
    D-03 & Recherche vectorielle native, avec repli applicatif
         & \emph{Écarté :} dépendance stricte à l'extension vectorielle. \emph{Critère :} la plateforme doit rester utilisable sans elle \\
    D-04 & Sélection du modèle salarial sur l'erreur de \emph{test}
         & \emph{Écarté :} sélection sur l'erreur d'entraînement. \emph{Critère :} un modèle qui mémorise se repère uniquement en test \\
    D-05 & Retrait d'une variable colinéaire du modèle
         & \emph{Écarté :} conservation de toutes les variables. \emph{Critère :} la colinéarité créait un mode de défaillance identifié \\
    D-06 & Assistant \emph{extractif} citant ses sources
         & \emph{Écarté :} génération libre. \emph{Critère :} une réponse non ancrée n'est pas vérifiable par le lecteur \\
    D-07 & Signature auto-hébergée avec empreinte et certificat
         & \emph{Écarté :} prestataire de signature externe. \emph{Critère :} maîtrise de la piste d'audit ; limite juridique assumée \\
    D-08 & Analyse de CV asynchrone, en session dédiée
         & \emph{Écarté :} analyse synchrone. \emph{Critère :} latence perçue indépendante de la taille du document \\
    \bottomrule
  \end{tabularx}
\end{table}

\section{Conclusion du chapitre}
\label{sec:arch_conclusion}

Ce chapitre a établi l'architecture de PIQBIT : un découpage en couches à
dépendance unidirectionnelle, un modèle de 25 tables piloté par migrations,
deux chemins d'entrée dont un public à jeton traité comme frontière de
confiance, et huit décisions d'architecture justifiées par leurs alternatives
écartées. Il a également déclaré les limites de la vue de déploiement :
absence de chiffrement du transport et d'environnement cible.

Le chapitre suivant décrit la réalisation de cette architecture, sous-système
par sous-système, en rapportant les difficultés rencontrées et la manière dont
elles ont été résolues.
